Instruction-tuned language models are updated one task at a time in practice, and each update erodes earlier skills.
Continual instruction-tuning benchmarks usually summarize this erosion with one end-of-sequence average under one or two task orders, which leaves the shape of forgetting, the role of task type, and the role of learning position unmeasured.
We test a three-part hypothesis: forgetting follows a power law, reasoning tasks are forgotten 3--5$\times$ faster than pattern-matching tasks, and early-learned tasks are forgotten disproportionately.
We fine-tune two small models (\llama and \qwen) with LoRA on six tasks in sequence, using a balanced Latin-square order design (18 sequential runs) and held-out loss measured every 25 steps.
None of the three parts holds as stated.
Forgetting is not uniform, but it is not a per-task power law either: it arrives as a fast, saturating hit at each task switch.
On \llama, a median 55\% of a 250-step stage's loss rise is present after 25 steps, and in a step-resolution experiment a saturating exponential beats a power law in 6 of 6 task pairs.
The reasoning/pattern forgetting ratio is 1.0--1.5 on \llama and 2.2--2.6 on \qwen for metric-based measures, with intervals that do not survive multiple-comparison correction.
Once the amount of subsequent training is fixed, early-learned tasks are not forgotten more on \llama.
The specific (old task, new task) pair is the strongest visible determinant of forgetting.
Our results suggest budgeting forgetting per task switch and per pair, not per training step or per task category.
